%  强凸强凹引理

\begin{proof}
For fixed $\blmbd$, the gradient of the reduced function is
\begin{equation}
    \nabla_{\bx}\mathcal D_\rho(\bx,\blmbd)
    =
    \nabla f(\bx)
    +J_{\calA}(\bx)^\top\bq(\bx,\blmbd).
    \label{eq:reduced-x-gradient}
\end{equation}
The map $u\mapsto u-\operatorname{prox}_{h/\rho}(u)$ is nonexpansive.
Consequently,
\begin{equation}
    \|\bq(\bx,\blmbd)-\bq(\bx',\blmbd)\|
    \le
    \rho G_{\calA}\|\bx-\bx'\|.
    \label{eq:q-lipschitz}
\end{equation}
Moreover, the proximal optimality condition gives
$\bq(\bx,\blmbd)\in\partial h(\by(\bx,\blmbd))$, and hence
$\|\bq(\bx,\blmbd)\|\le L_h$.  Assumptions~\ref{ass:smoothness}
and~\ref{ass:h} therefore imply
\begin{equation}
\begin{aligned}
    &\|\nabla_{\bx}\mathcal D_\rho(\bx,\blmbd)
       -\nabla_{\bx}\mathcal D_\rho(\bx',\blmbd)\| \\
    &\quad\le
    \bigl(L_f+L_{\calA}L_h+\rho G_{\calA}^2\bigr)
    \|\bx-\bx'\|
    =L_{\calD,\bx}\|\bx-\bx'\|.
\end{aligned}
\label{eq:reduced-gradient-lipschitz}
\end{equation}
Thus $\mathcal D_\rho(\cdot,\blmbd)$ is
$L_{\calD,\bx}$-weakly convex.  Adding the term
$\|\bx-\bz\|^2/(2\gamma)$ shows that
$\mathcal P(\cdot,\blmbd;\bz,\bmu)$ is
$(1/\gamma-L_{\calD,\bx})$-strongly convex.  Since
$1/\gamma>2L_{\calD,\bx}$, its strong-convexity modulus is at least
$L_{\calD,\bx}$.

For fixed $\bx$, the function $\mathcal D_\rho(\bx,\cdot)$ is concave,
because it is the pointwise infimum over $\by$ of functions affine in
$\blmbd$.  Adding
$-\|\blmbd-\bmu\|^2/(2s)$ makes
$\mathcal P(\bx,\cdot;\bz,\bmu)$ $1/s$-strongly concave.
\end{proof}


%  Basic estimate of the reduced potential
\begin{proof}
The \(\bx\)-update is a gradient step on \(\mathcal P_k\) with
\((\blmbd^k,\bz^k,\bmu^k)\) fixed.  The descent lemma and
\(\Delta_{\bx}^{k+1}=-\eta_k\nabla_{\bx}\mathcal P_k(\bx^k,\blmbd^k)\)
give
\begin{equation}
    \mathcal P_k(\bx^{k+1},\blmbd^k)
    -\mathcal P_k(\bx^k,\blmbd^k)
    \le
    -\left(\frac{1}{\eta_k}-\frac{L_x}{2}\right)
    \|\Delta_{\bx}^{k+1}\|^2.
    \label{eq:x-descent}
\end{equation}

With \(\bx^{k+1}\) fixed, the dual update satisfies
\(\Delta_{\blmbd}^{k+1}=\tau_k
\nabla_{\blmbd}\mathcal P_k(\bx^{k+1},\blmbd^k)\).  Applying the smoothness
bound \eqref{eq:dual-gradient-smoothness} yields
\begin{equation}
    \mathcal P_k(\bx^{k+1},\blmbd^{k+1})
    -\mathcal P_k(\bx^{k+1},\blmbd^k)
    \le
    \left(\frac{1}{\tau_k}+\frac{L_{\lambda,k}}{2}\right)
    \|\Delta_{\blmbd}^{k+1}\|^2.
    \label{eq:auxiliary-dual-ascent}
\end{equation}

The center update gives the exact identity
\begin{equation}
\begin{aligned}
    &\mathcal P(\bx^{k+1},\blmbd^{k+1};\bz^{k+1},\bmu^k)
    -\mathcal P(\bx^{k+1},\blmbd^{k+1};\bz^k,\bmu^k) \\
    &\qquad=
    -\frac{\alpha_k(2-\alpha_k)}{2\gamma}
    \|\bx^{k+1}-\bz^k\|^2.
\end{aligned}
\label{eq:center-drift}
\end{equation}

Finally, the multiplier-center update gives
\begin{equation}
\begin{aligned}
    &-\frac{1}{2s_k}\|\blmbd^{k+1}-\bmu^{k+1}\|^2
    +\frac{1}{2s_k}\|\blmbd^{k+1}-\bmu^k\|^2 \\
    &\qquad=
    \frac{\beta_k(2-\beta_k)}{2s_k}
    \|\blmbd^{k+1}-\bmu^k\|^2.
\end{aligned}
\label{eq:mu-drift}
\end{equation}
Changing \(s_k\) to \(s_{k+1}\) contributes the final term in
\eqref{eq:basic-one-step}.  Adding
\eqref{eq:x-descent}, \eqref{eq:auxiliary-dual-ascent},
\eqref{eq:center-drift}, and the two parameter-drift identities proves
\eqref{eq:basic-one-step}.
\end{proof}
